\scanpage{38}
% Pencil archive foliation: 37.
Si $\Sigma^+(X)=\text{ens des classes d'objets simples}$
de $\mathcal{M}^+(X)$,
\[
\Sigma(X)=\underset{\substack{\longrightarrow\\
\text{par }e\unsure{\rightsquigarrow|\mapsto}\xi e}}{\lim}\Sigma^+(X),
\]
on aura par \deleted{$\mathcal{M}(K)$} un diagramme comm.
\diagram{\begin{tikzcd}
\mathrm{M}^+(X) \arrow[d,hook] \arrow[r,equal] &
\mathbb{Z}^{(\Sigma^+(X))} \arrow[d,hook]\\
\mathrm{M}(X) \arrow[r,equal] &
\mathbb{Z}^{(\Sigma(X))}
\end{tikzcd}}
Alors $\mathcal{M}_i^+(X)$ est défini en termes
d'une partie $\Sigma_{(i)}^+(X)$ de
$\unsure{\Sigma_i(X)|\Sigma(X)}$,
d'où
\[
\mathrm{M}_i^+(X)=\mathbb{Z}^{\Sigma_{(i)}^+(X)}
\]
ce qui donne une filtration croissante
sur l'anneau $\mathrm{M}^+(X)$, et une
filtration analogue sur $\mathrm{M}(X)$.
En fait, posant
\[
\Sigma_i^+(X)=\Sigma_{(i)}^+(X)-\Sigma_{(i-1)}^+(X)
\]
on aura \unsure{$=$|$\overline{\mathrm m}$}
\[
\Sigma_i^+(X).\Sigma_j^+(X)\subset
\mathbb{Z}^{\Sigma_{i+j}^+(X)}
\]
donc $\mathrm{M}^+(X)$ devient un \uline{anneau}
